% GENERATED FILE. Edit canonical lessons, then run npm run generate:books.
% canonical-source-hash: fnv1a64:c955a381eb2c2436
% canonical-lessons: TE-C190-karentu, TE-C190-yantram, TE-C190-sopha, TE-C190-uyyala, TE-C190-tivaci

\chapter{Electricity, Machine, Sofa, Swing, Carpet}
\label{ch:te-electricity-machine-sofa-swing-carpet}
\hlchaptermodality{190}

\begin{chapteropening}
\textbf{By the end of this chapter:} \emph{I can say electricity, a machine, a sofa, a swing and a carpet.}

Complete the last lesson of chapter 190: I can say electricity, a machine, a sofa, a swing and a carpet.
\end{chapteropening}

\section[\texorpdfstring{\te{కరెంటు} (kareṇṭu)}{కరెంటు (kareṇṭu)}]{\te{కరెంటు} (kareṇṭu) — electricity, power}

\label{lesson:TE-C190-karentu}



\begin{quote}
Before the new one: say the Telugu for a hobby, then the Telugu for a novel.
\end{quote}

\subsection*{You'll want to know: \te{కరెంటు}}
\textbf{\te{కరెంటు}} — \emph{kareṇṭu} — \textquotedblleft{}electricity, power\textquotedblright{}.

\begin{grammarlens}[title={What you've built}]
One more word for this chapter.
\end{grammarlens}

\subsection*{Your turn}
\begin{itemize}
\raggedright
  \item \emph{Say it:} \emph{kareṇṭu}
  \item \emph{Say it:} \emph{kareṇṭu}, once more
  \item \emph{Say it:} say \emph{kareṇṭu}
  \item \emph{Recall:} say the Telugu for a hobby, then the Telugu for a novel, then say \emph{kareṇṭu} again
\end{itemize}

\begin{tcolorbox}[breakable,colback=teal!4,colframe=teal!35!black,title={Before you move on}]
What does \te{కరెంటు} mean? (\textquotedblleft{}Electricity, power\textquotedblright{}.) Say it once more.
\end{tcolorbox}
\section[\texorpdfstring{\te{యంత్రం} (yantraṁ)}{యంత్రం (yantraṁ)}]{\te{యంత్రం} (yantraṁ) — a machine}

\label{lesson:TE-C190-yantram}



\begin{quote}
Before the new one: say the Telugu for a novel, then the Telugu for electricity.
\end{quote}

\subsection*{You'll want to know: \te{యంత్రం}}
\textbf{\te{యంత్రం}} — \emph{yantraṁ} — \textquotedblleft{}a machine\textquotedblright{}.

\begin{grammarlens}[title={What you've built}]
One more word for this chapter.
\end{grammarlens}

\subsection*{Your turn}
\begin{itemize}
\raggedright
  \item \emph{Say it:} \emph{yantraṁ}
  \item \emph{Say it:} \emph{yantraṁ}, once more
  \item \emph{Say it:} say \emph{yantraṁ}
  \item \emph{Recall:} say the Telugu for a novel, then the Telugu for electricity, then say \emph{yantraṁ} again
\end{itemize}

\begin{tcolorbox}[breakable,colback=teal!4,colframe=teal!35!black,title={Before you move on}]
What does \te{యంత్రం} mean? (\textquotedblleft{}A machine\textquotedblright{}.) Say it once more.
\end{tcolorbox}
\section[\texorpdfstring{\te{సోఫా} (sōphā)}{సోఫా (sōphā)}]{\te{సోఫా} (sōphā) — a sofa}

\label{lesson:TE-C190-sopha}



\begin{quote}
Before the new one: say the Telugu for electricity, then the Telugu for a machine.
\end{quote}

\subsection*{You'll want to know: \te{సోఫా}}
\textbf{\te{సోఫా}} — \emph{sōphā} — \textquotedblleft{}a sofa\textquotedblright{}.

\begin{grammarlens}[title={What you've built}]
One more word for this chapter.
\end{grammarlens}

\subsection*{Your turn}
\begin{itemize}
\raggedright
  \item \emph{Say it:} \emph{sōphā}
  \item \emph{Say it:} \emph{sōphā}, once more
  \item \emph{Say it:} say \emph{sōphā}
  \item \emph{Recall:} say the Telugu for electricity, then the Telugu for a machine, then say \emph{sōphā} again
\end{itemize}

\begin{tcolorbox}[breakable,colback=teal!4,colframe=teal!35!black,title={Before you move on}]
What does \te{సోఫా} mean? (\textquotedblleft{}A sofa\textquotedblright{}.) Say it once more.
\end{tcolorbox}
\section[\texorpdfstring{\te{ఉయ్యాల} (uyyāla)}{ఉయ్యాల (uyyāla)}]{\te{ఉయ్యాల} (uyyāla) — a swing; a cradle}

\label{lesson:TE-C190-uyyala}



\begin{quote}
Before the new one: say the Telugu for a machine, then the Telugu for a sofa.
\end{quote}

\subsection*{You'll want to know: \te{ఉయ్యాల}}
\textbf{\te{ఉయ్యాల}} — \emph{uyyāla} — \textquotedblleft{}a swing; a cradle\textquotedblright{}.

\begin{grammarlens}[title={What you've built}]
One more word for this chapter.
\end{grammarlens}

\subsection*{Your turn}
\begin{itemize}
\raggedright
  \item \emph{Say it:} \emph{uyyāla}
  \item \emph{Say it:} \emph{uyyāla}, once more
  \item \emph{Say it:} say \emph{uyyāla}
  \item \emph{Recall:} say the Telugu for a machine, then the Telugu for a sofa, then say \emph{uyyāla} again
\end{itemize}

\begin{tcolorbox}[breakable,colback=teal!4,colframe=teal!35!black,title={Before you move on}]
What does \te{ఉయ్యాల} mean? (\textquotedblleft{}A swing; a cradle\textquotedblright{}.) Say it once more.
\end{tcolorbox}
\section[\texorpdfstring{\te{తివాచీ} (tivācī)}{తివాచీ (tivācī)}]{\te{తివాచీ} (tivācī) — a carpet, a rug}

\label{lesson:TE-C190-tivaci}



\begin{quote}
Before the new one: say the Telugu for a sofa, then the Telugu for a swing.
\end{quote}

\subsection*{You'll want to know: \te{తివాచీ}}
\textbf{\te{తివాచీ}} — \emph{tivācī} — \textquotedblleft{}a carpet, a rug\textquotedblright{}.

\begin{grammarlens}[title={What you've built}]
One more word for this chapter.
\end{grammarlens}

\subsection*{Your turn}
\begin{itemize}
\raggedright
  \item \emph{Say it:} \emph{tivācī}
  \item \emph{Say it:} \emph{tivācī}, once more
  \item \emph{Say it:} say \emph{tivācī}
  \item \emph{Recall:} say the Telugu for a sofa, then the Telugu for a swing, then say \emph{tivācī} again
\end{itemize}

\begin{tcolorbox}[breakable,colback=teal!4,colframe=teal!35!black,title={Before you move on}]
What does \te{తివాచీ} mean? (\textquotedblleft{}A carpet, a rug\textquotedblright{}.) Say it once more.
\end{tcolorbox}
